\documentclass[11pt, a4paper]{article}

% --- Packages ---
\usepackage[margin=1in]{geometry} % Standard margins
\usepackage[utf8]{inputenc}
\usepackage[T1]{fontenc}
\usepackage{lmodern} % Sharp, scalable fonts
\usepackage{amsmath, amssymb} % Math symbols
\usepackage{xcolor} % Colors
\usepackage{tcolorbox} % Beautiful boxes
\tcbuselibrary{skins, breakable}
\usepackage{enumitem} % Better lists
\usepackage{fancyhdr} % Headers/Footers
\usepackage{hyperref} % Clickable links
\usepackage{titlesec} % Custom section styling

% --- Color Definitions ---
\definecolor{techblue}{RGB}{0, 51, 102}
\definecolor{alertred}{RGB}{153, 0, 0}
\definecolor{softgray}{RGB}{245, 245, 245}
\definecolor{codegreen}{RGB}{0, 100, 0}

% --- Custom Box Styles ---

% 1. The "Mental Model" Box (Blue)
\newtcolorbox{mentalmodel}[1]{
  colback=techblue!5!white,
  colframe=techblue,
  fonttitle=\bfseries,
  title={#1},
  enhanced,
  boxrule=0.5mm,
  attach boxed title to top left={xshift=0.5cm, yshift=-2mm},
  boxed title style={colback=techblue, sharp corners},
  breakable
}

% 2. The "Equation" Box (Gray)
\newtcolorbox{eqbox}[1]{
  colback=white,
  colframe=gray!50!black,
  fonttitle=\bfseries,
  title={#1},
  enhanced,
  attach boxed title to top center={yshift=-2mm},
  boxed title style={colback=gray!50!black},
  sharp corners,
  breakable
}

% 3. The "Warning/Crucial" Box (Red)
\newtcolorbox{crucial}[1]{
  colback=alertred!5!white,
  colframe=alertred,
  fonttitle=\bfseries,
  title={#1},
  enhanced,
  boxrule=0.5mm,
  attach boxed title to top left={xshift=0.5cm, yshift=-2mm},
  boxed title style={colback=alertred},
  breakable
}

% --- Page Layout ---
\pagestyle{fancy}
\fancyhf{}
\lhead{\textbf{Static Timing Analysis}}
\rhead{Teaching Notes}
\cfoot{\thepage}
\renewcommand{\headrulewidth}{1pt}

% --- Title Section ---
\title{\textbf{\Huge Static Timing Analysis}\\ \Large The Physics of Digital Time}
\author{\textit{Structured Learning Notes}}
\date{\today}

\begin{document}

\maketitle
\thispagestyle{empty}
\tableofcontents
\newpage

\section{The Core Problem: Physics vs. Logic}

Forget software tools for a moment. Think physics. A digital circuit is essentially a race. Signals leave one register, travel through logic gates and wires, and must arrive at the next register at the exact right moment.

\begin{itemize}
    \item \textbf{Too Late:} Setup failure $\rightarrow$ The wrong value is captured.
    \item \textbf{Too Early:} Hold failure $\rightarrow$ The new value overwrites the old one too fast.
\end{itemize}

\begin{mentalmodel}{STA's One Question}
``For every possible path in this design, under worst-case conditions, will data arrive at the right time?''
\end{mentalmodel}

\section{Why Simulation Fails (And STA Exists)}

\subsection{The Simulation Mindset}
Simulation is like watching a movie. You choose input vectors, watch what happens, and hope you covered all cases.
\begin{itemize}
    \item \textbf{Problems:} It misses corner cases, it is non-exhaustive, and it is computationally impossible for million-path designs.
\end{itemize}

\subsection{The STA Mindset}
STA is like a mathematical proof.
\begin{itemize}
    \item No input vectors required.
    \item No guessing.
    \item \textbf{Topological Propagation:} It pushes timing through gates, accumulates delay, and stops at endpoints.
\end{itemize}

\section{Vector-less Analysis: The Mechanism}

STA does not care about logical values (``0'' or ``1''). It ignores functional behavior and focuses entirely on \textbf{delay math}.

\begin{description}
    \item[Pipes:] Gates + Wires
    \item[Flow Time:] Delay
    \item[Reservoirs:] Registers
\end{description}

\section{Timing Goals \& Frequency}

Every synchronous chip lives under one law: \textbf{Everything must fit inside one clock period.}

If your clock is 1 GHz ($T_{period} = 1\text{ns}$):
\begin{itemize}
    \item All logic between registers must finish in $< 1\text{ns}$ (minus uncertainty margins).
    \item STA verifies frequency targets, signal collisions, and data integrity.
\end{itemize}

\section{Variability: Why Delay isn't a Single Number}

Delays change constantly due to physics. STA uses \textbf{PVT} (Process, Voltage, Temperature) variations.
\begin{enumerate}
    \item \textbf{Process:} Transistors vary during manufacturing.
    \item \textbf{Voltage:} IR drop slows down switching.
    \item \textbf{Temperature:} Hot silicon is generally slower.
    \item \textbf{Noise:} Crosstalk and jitter.
\end{enumerate}

\begin{crucial}{Min/Max Analysis}
STA must always assume:
\begin{itemize}
    \item \textbf{Worst Case Slow:} For Setup checks (Late Analysis).
    \item \textbf{Worst Case Fast:} For Hold checks (Early Analysis).
\end{itemize}
\end{crucial}

\section{Registers: The Traffic Lights}

Registers define the timing windows. They are not just storage; they are constraints.

\begin{eqbox}{Setup Constraint (Max Delay)}
Data must arrive \textbf{before} the clock edge.
\[ T_{data\_delay} + T_{setup} < T_{clock\_period} \]
\textit{Violation:} Data arrives too late.
\end{eqbox}

\begin{eqbox}{Hold Constraint (Min Delay)}
Data must stay stable \textbf{after} the clock edge.
\[ T_{data\_delay} > T_{hold} \]
\textit{Violation:} Data changes too early (racing through).
\end{eqbox}

\section{The Anatomy of a Timing Path}

STA enumerates every possible path between valid \textbf{Startpoints} and \textbf{Endpoints}.

\begin{center}
\begin{tabular}{|c|c|}
\hline
\textbf{Startpoints} & \textbf{Endpoints} \\
\hline
Input Ports & Output Ports \\
Clock pins of Flip-Flops & Data pins of Flip-Flops \\
\hline
\end{tabular}
\end{center}

\section{Special Path Types}

\subsection{False Paths}
Some paths exist physically but not logically (e.g., mutually exclusive MUX inputs). STA doesn't know logic intent. You must constrain these as \texttt{set\_false\_path}, otherwise STA will report violations for things that never happen.

\subsection{Multi-Cycle Paths (MCP)}
Logic intentionally designed to take $>1$ cycle (e.g., complex arithmetic). By default, STA assumes everything is single-cycle. MCP constraints relax the \textbf{Required Time}.

\section{The Mathematics of Timing}

\subsection{Required Time vs. Arrival Time}
\begin{itemize}
    \item \textbf{Required Time:} The deadline. ``By when \textit{must} data arrive?''
    \item \textbf{Arrival Time:} The reality. ``When \textit{does} data actually get there?''
\end{itemize}

\subsection{Slack}
Slack is the single most important number in a timing report. It represents your \textbf{Margin}.

\begin{eqbox}{The Slack Equation}
\[ \text{Slack} = \text{Required Time} - \text{Arrival Time} \]
\end{eqbox}

\begin{itemize}
    \item \textbf{Positive Slack:} Safe.
    \item \textbf{Zero Slack:} Razor's edge (exact met timing).
    \item \textbf{Negative Slack:} \textcolor{alertred}{\textbf{FAILURE}}.
\end{itemize}

The path with the most negative slack is your \textbf{Critical Path}.

\section{Early vs. Late Analysis (Crucial)}

\begin{mentalmodel}{Two Universes of Analysis}
\begin{itemize}
    \item \textbf{Late Analysis (Setup):} Assumes everything is \textbf{slow}. Checks if the signal arrives \textit{before} the next clock edge.
    \item \textbf{Early Analysis (Hold):} Assumes everything is \textbf{fast}. Checks if the signal stays stable \textit{long enough} after the current clock edge.
\end{itemize}
\end{mentalmodel}

\section{Summary: The Mental Model}

Think of Static Timing Analysis as a physics-based scheduling problem with strict deadlines (Required Time) and actual physical estimates (Arrival Time).

\begin{itemize}
    \item STA is not pessimistic by accident; it is pessimistic by \textbf{design}.
    \item If you pass STA, you are mathematically guaranteed to work on silicon (within the modeled physics).
\end{itemize}

\end{document}